\section{Bloque V: catálogo actualizable y operación integrada}
\label{sec:bloque5}
\textbf{Fecha: 18 de octubre de 2026.} Implementación local de P17--P20.
La evidencia se conserva en \code{validation/block5/}. Las pruebas con adaptadores
de control verifican contratos y fallos, no relevancia musical ni descargas reales.

\subsection{P17: generaciones coherentes del catálogo}
Antes se conservaban caches independientes del CSV, mapa e índices. Reemplazar
archivos mientras una búsqueda estaba en curso podía combinar generaciones.
Ahora \code{catalog.py} publica los cuatro artefactos en un directorio inmutable,
valida dimensiones, alineación, normalización y hashes SHA-256 y cambia el puntero
actual mediante una transacción SQLite. El historial publicado distingue las
generaciones válidas de los directorios huérfanos de una publicación fallida.
El primer arranque copia los archivos heredados de \code{output/}, sin modificarlos.

Cada solicitud web o búsqueda CLI fija un snapshot durante todo su recorrido.
Una consulta iniciada antes de una actualización termina con su generación; la
siguiente adopta la nueva. Los modelos se reutilizan, pero los recursos de búsqueda
se superponen desde el contexto de la solicitud. No se agregan caches globales de
resultados. Los helpers heredados sin contexto conservan su comportamiento estático.
Las vistas de metadatos son copias y los índices expuestos no permiten agregar ni
borrar vectores. Ante un nuevo puntero inválido, un servidor que ya posee un snapshot
continúa con él e informa una advertencia; un arranque sin catálogo válido falla.

Los publicadores se serializan: agregar canciones recarga la última generación
dentro de la transacción. Reconstruir todo exige que la versión inicial no haya
cambiado; así una ingesta concurrente no desaparece silenciosamente. Los enlaces
de reproducción incluyen versión e ID, evitando que un ID reutilizado reproduzca
otra canción. Se comprueba que la versión pertenezca al historial publicado.
El coste es copiar índices y catálogo en cada publicación: es apropiado para el
catálogo pequeño actual, pero debe medirse antes de escalar. No se garantiza
durabilidad frente a pérdida de energía ni se implementa todavía rollback operativo.

\subsection{P18: decisión sobre persistencia y progreso}
No se recuperan checkpoints del grafo de búsqueda: sus juicios y recursos podrían
pertenecer a otra generación, y no existe evidencia de que el ahorro compense esa
complejidad. La búsqueda sigue siendo independiente por solicitud. Para ingesta,
sí se persisten trabajos en SQLite: estado, etapa, progreso, metadatos, propietario,
lease, resultado y error. La interfaz consulta el estado cada 1,5 segundos y conserva
el último ID en almacenamiento local. Recargar consulta ese ID; no vuelve a enviar
la canción ni reanuda automáticamente la búsqueda.

Un trabajo pendiente puede ser reclamado una sola vez entre procesos. Un trabajo
en ejecución pierde su propiedad al vencer el lease de 60 segundos y queda
interrumpido. Un trabajador vencido no puede renovar ni publicar. El usuario puede
enviar de nuevo el enlace; la publicación es idempotente por trabajo y URL. Si hubo
una caída tras publicar y antes de registrar el resultado, el nuevo envío detecta
la canción existente. La cola admite hasta 32 trabajos activos y deduplica enlaces.
Se eligió polling por sencillez y recuperación por ID; SSE queda diferido hasta
demostrar una necesidad de menor latencia o menor tráfico.

\subsection{P19: ingesta opcional sin bloquear la búsqueda}
La función permanece desactivada por defecto mediante
\code{SPELLFIND_INGEST_ENABLED=false}. El adaptador nativo usa yt-dlp para MP3 con
FFmpeg, prioriza subtítulos oficiales y utiliza Whisper CPU como alternativa.
Un fallo al descargar subtítulos no descarta el audio ya descargado. Se aceptan
enlaces HTTPS de vídeos individuales de YouTube, con duración conocida de hasta
15 minutos, límites de tamaño y tiempos de espera. Se rechazan directos y enlaces
arbitrarios. Estos límites no equivalen a cuotas globales de almacenamiento.

Las letras se codifican con prefijo \code{search_document:}, fragmentos de 2000
caracteres y media normalizada de vectores de dimensión 256. CLAP procesa ventanas
de 10 segundos al inicio, centro y final a 48 kHz y produce una media normalizada
de dimensión 512. La cobertura es parcial y se registra. La receta del catálogo
importado no demuestra igualdad de revisiones ni de agregación con los vectores
nuevos; esa compatibilidad debe evaluarse antes de afirmar calidad equivalente.

Sólo se publica después de completar metadatos y ambos vectores. No se inventan
letras vacías ni se usa la fecha de subida como año de lanzamiento. Artista, género
y año pueden ser aportados explícitamente; la procedencia queda registrada. Una
canción sin transcripción utilizable falla, incluso si es instrumental: el esquema
alineado actual requiere ambas modalidades. Esto es una limitación pendiente,
no una mejora sobre recuperación por audio de instrumentales.

La API añade \code{GET /ingest/status}, \code{POST /ingest} y
\code{GET /ingest/jobs/<id>}. Responde 202 al encolar, 400 ante datos inválidos,
403 ante origen cruzado, 429 ante saturación y 503 si la capacidad está desactivada
o incompleta. La interfaz muestra etapas, errores y reproducción versionada.
El servicio está concebido para uso local; la comprobación de origen no sustituye
autenticación para exposición pública. La extensión
\code{requirements/ingestion.in} es experimental y no tiene un lock verificado.
En este entorno faltan yt-dlp, FFmpeg y Whisper: no se realizó una ingesta real.
Referencias técnicas: \url{https://github.com/yt-dlp/yt-dlp},
\url{https://github.com/openai/whisper} y \url{https://github.com/LAION-AI/CLAP}.

\subsection{P20: comprobaciones y límites del cierre}
Las regresiones ejercitan publicación fallida, hashes corruptos, snapshots
inmutables, dos publicadores, comparación de versión, búsqueda en curso durante
actualización, propiedad y vencimiento de trabajos, errores del adaptador,
subtítulos y fallback, contratos HTTP, audio con rangos y consistencia web/CLI.
Las evaluaciones fijan un snapshot para toda la serie y registran su manifest;
la comparación rechaza catálogos diferentes. Los resultados ejecutados y sus
comandos quedan en los informes del bloque. No se confunden adaptadores de
control con ejecución de proveedores ni una comprobación estática con compilación PDF.

\subsection{Pendientes adicionales y prioridad}
\begin{description}
\item[E17, P17/P20: retención y recuperación.] Definir cuotas, limpieza segura de
medios, trabajos y generaciones huérfanas, backups y restauración. Actualmente se
retienen para no eliminar recursos de consultas o enlaces todavía activos.
\item[E18, P01/P19/P20: entorno e integración reales.] Fijar dependencias opcionales
en entorno separado; verificar FFmpeg, subtítulos, Whisper, revisiones de Nomic/CLAP
y compatibilidad con vectores históricos. Después medir relevancia humana,
latencia y memoria reales. Mantener pendientes P04, P10, P12 y P16 donde corresponda.
\item[E19, P06/P19: modalidades opcionales.] Diseñar disponibilidad por canción e
índices con mapas independientes para admitir instrumentales sin letras ficticias.
Revisar normalización y explicación de pesos antes de implementarlo.
\end{description}
Primero se valida E18 en un entorno reproducible; E17 precede al uso sostenido y
E19 a anunciar ingesta general de instrumentales. Ninguno se considera cerrado
por haber añadido la infraestructura de este bloque.
